\subsection{Measuring the Migration Cost}
\label{sec:online:M}

Everything in \S\ref{sec:online} turns on $M$, the cost of moving a primary's
state to a different backup relative to one round of ordinary replication. Up to
this point $M$ was a swept parameter, for a reason worth stating plainly: DEFT's
backup does not keep the state it is sent. Its \texttt{update\_state} handler
prints the payload and returns, with the one line that would have persisted it
commented out, so the backup holds nothing and could not in fact take over. The
replication experiments in \S\ref{sec:results} consequently ship a single
hard-coded dictionary entry of about thirty bytes. That is sufficient for
measuring the cost of an RPC round trip, which is what those experiments set out
to do, but it means there is no state whose transfer could be timed.

We therefore implement the missing component: a connection-tracking flow table of
the kind a stateful NAT maintains, keyed by the 5-tuple in the same
comma-separated form DEFT's placeholder already uses, with an entry holding
packet and byte counters, a translation, a last-seen timestamp and a TCP state.
It supports the three operations replication actually requires---fold an
incremental update into the table, serialize the whole table, and load it---so
that a backup is a warm standby rather than a sink. The table is filled by a
flow-arrival workload with Poisson arrivals and heavy-tailed, truncated-Pareto
lifetimes, with idle eviction. Steady-state occupancy agrees with the Little's-law
prediction to within 10\%, which we use as a check on the generator rather than as
a result. The workload parameters are not measured and are swept over a
$16\times$ range of arrival rates for that reason; replacing the generator with a
packet trace changes nothing downstream, since every later stage consumes only the
resulting table.

\paragraph{The encoding DEFT uses costs more than the state does.}
Table~\ref{tab:encoding} reports the serialized size of the same table under
three encodings. A fixed-width binary layout needs 44 bytes per flow. DEFT's
XML-RPC marshalling needs over seven times that. Part of the gap is per-value
markup, but part of it is a hard limitation: XML-RPC's \texttt{<int>} is signed
32-bit, so an IPv4 address from 128.0.0.0 upward does not fit and a byte counter
overflows past 2\,GiB. Any real NAT state must therefore cross this transport as
text. This compounds the transport finding of \S\ref{sec:p99}: the shipped
RPC stack distorts the latency measurement \emph{and} inflates the quantity that
governs whether adaptation is worthwhile.

\paragraph{The migration cost is tens of rounds, not a few.}
Table~\ref{tab:mrounds} gives $M$ in rounds. Across the sweep it lands between
roughly 23 and 41, where the online sweeps had assumed $3$. The reason it is
large is structural rather than incidental: a round replicates only the flows that
changed during that round---a few percent of the table, per
Table~\ref{tab:occupancy}---whereas a migration ships all of it. $M$ is
consequently close to the reciprocal of the delta share, which makes it a
property of traffic churn rather than of the wire format, since a costlier
encoding inflates the full transfer and the per-round delta in equal measure.
The encoding does move the cost in absolute time, and substantially. At 10\,Gb/s,
migrating a 62{,}000-flow table takes about 165\,ms packed and about 2.1\,s over
DEFT's XML-RPC; at the top of the sweep, a 248{,}000-flow table takes 8.8\,s over
XML-RPC against 0.7\,s packed.

\paragraph{What this does to the conclusions.}
Re-running the sweeps of \S\ref{sec:online:results} at the measured
$M$ sharpens them in one direction and overturns one of them.

\textsc{snapshot} becomes far worse: its excess ratio at the gentlest churn rises
from $11.2$ to $46.1$, and it now spends $96\%$ of its total cost on migration.
\textsc{periodic} rebalancing, the sensible-looking operational compromise, spends
$60\%$ of its cost on migration at low churn and is no longer defensible either.
The cost of acting on stale or transient information scales with $M$, and $M$ was
underestimated by a factor of three to five.

The overturned result is the comparison between hysteresis and doing nothing. At
$M=3$ the two traded places across the persistence sweep, and we reported that no
regime clearly favoured reacting to degradation. At the measured $M$ there is
essentially no trading. At the top of the measured range ($M = 13.7$)
\textsc{static} is ahead at every degradation duration we test, from $2$ to $300$
rounds, at $1.15$--$1.20$ against $1.22$--$1.31$. At the bottom ($M = 7.7$) it is
ahead at four of the six durations, level at one, and behind at one by $0.01$.
The largest advantage hysteresis holds anywhere in the measured range is that
$0.01$, against deficits of up to $0.12$. With a realistic
migration cost, reacting to degradation does not pay anywhere in this model's
measured range; only reacting to hard failure does.

The measured crossover says when that would change. A severe degradation---the
60\,ms of the \texttt{link\_congested} kind---must persist for about $5.5$
seconds before migrating away from it repays the round trip, under a compact
encoding at 10\,Gb/s. Under DEFT's own encoding the same degradation must last
$71$ seconds, and a mild 5\,ms one over fourteen minutes. These are the numbers an operator can
check against their own failure durations, and they are the practical output of
the online analysis: not an algorithm to deploy, but a condition to evaluate
before deciding whether any algorithm is needed.
